%Author: Marco Miani
%SWAN (developed by SWAN group, TU Delft, The Netherlands) is a wave spectral numerical model.
%For Simlating WAves Nearshore, it is necessary to define spatial grids of
%physical dominant factors (wind friction, dissipation) as well as define a COMPUTATIONAL
%grid on which the model performs its (spectral) calculations: budgeting energy spectra over
%each cell of the (computational) grid. Grids might have different spatial resolution and extension.

% Minimalist adaptation: shared roles, native Latex heads, and revised spacing.
\documentclass[tikz,border=10bp]{standalone}
\usepackage{minimalist-profile}
\begin{document}



\begin{tikzpicture}[ms base,on grid]

    %slanting: production of a set of n 'laminae' to be piled up. N=number of grids.
    \begin{scope}[
            yshift=-83,every node/.append style={
            yslant=0.5,xslant=-1},yslant=0.5,xslant=-1
            ]
        % opacity to prevent graphical interference
        \fill[msPaper,fill opacity=1] (0,0) rectangle (5,5);
        \draw[step=4mm,ms context,msColor5] (0,0) grid (5,5); %defining grids
        \draw[step=1mm, msColor1,ms context] (3,1) grid (4,2);  %Nested Grid
        \draw[msColor5,ms structure] (0,0) rectangle (5,5);%marking borders
        \fill[msColor1] (0.05,0.05) rectangle (0.35,0.35);
        %Idem as above, for the n-th grid:
    \end{scope}

    \begin{scope}[
        yshift=0,every node/.append style={
            yslant=0.5,xslant=-1},yslant=0.5,xslant=-1
                     ]
        \fill[msPaper,fill opacity=1] (0,0) rectangle (5,5);
        \draw[msColor5,ms structure] (0,0) rectangle (5,5);
        \draw[step=5mm,ms context,msColor5] (0,0) grid (5,5);
    \end{scope}

    \begin{scope}[
        yshift=90,every node/.append style={
        yslant=0.5,xslant=-1},yslant=0.5,xslant=-1
                     ]
        \fill[msPaper,fill opacity=1] (0,0) rectangle (5,5);
        \draw[step=10mm,ms context,msColor5] (1,1) grid (4,4);
        \draw[msColor5,ms structure] (1,1) rectangle (4,4);
        \draw[msColor5,ms dashed] (0,0) rectangle (5,5);
    \end{scope}

    \begin{scope}[
        yshift=170,every node/.append style={
            yslant=0.5,xslant=-1},yslant=0.5,xslant=-1
          ]
        \fill[msPaper,fill opacity=1] (0,0) rectangle (5,5);
        \draw[step=10mm,ms context,msColor5] (2,2) grid (5,5);
        \draw[step=2mm, msColor2] (2,2) grid (3,3);
        \draw[msColor5,ms structure] (2,2) rectangle (5,5);
        \draw[msColor5,ms dashed] (0,0) rectangle (5,5);
    \end{scope}

    \begin{scope}[
        yshift=-170,every node/.append style={
        yslant=0.5,xslant=-1},yslant=0.5,xslant=-1
                  ]
        %marking border
        \draw[msColor5,ms structure] (0,0) rectangle (5,5);

        %drawing corners (P1,P2, P3): only 3 points needed to define a plane.
        \draw [fill=msColor2](0,0) circle[radius={0.5*\msMarkerEmphasis}] ;
        \draw [fill=msColor2](0,5) circle[radius={0.5*\msMarkerEmphasis}];
        \draw [fill=msColor2](5,0) circle[radius={0.5*\msMarkerEmphasis}];
        \draw [fill=msColor2](5,5) circle[radius={0.5*\msMarkerEmphasis}];

        %drawing bathymetric hypotetic countours on the bottom grid:
        \draw [ms emphasis,msColor5](0,1) parabola bend (2,2) (5,1)  ;
        \draw [ms dashed] (0,1.5) parabola bend (2.5,2.5) (5,1.5) ;
        \draw [ms dashed] (0,2) parabola bend (2.7,2.7) (5,2)  ;
        \draw [ms dashed] (0,2.5) parabola bend (3.5,3.5) (5,2.5)  ;
        \draw [ms dashed] (0,3.5)  parabola bend (2.75,4.5) (5,3.5);
        \draw [ms dashed] (0,4)  parabola bend (2.75,4.8) (5,4);
        \draw [ms dashed] (0,3)  parabola bend (2.75,3.8) (5,3);
        \coordinate (shore-label) at (2.8,1);
        \coordinate (shore-target) at (2,1.99);
    \end{scope} %end of drawing grids

    % Labels stay outside the slanted geometry so glyphs retain their physical shape.
    \draw[ms arrow,draw=msColor5] (shore-label) node[right,text=msColor5] {Shoreline}
      to[out=180,in=270] (shore-target);
    %putting arrows and labels:
    \draw[ms arrow] (6.2,2) node[right]{Bathymetry}
         to[out=180,in=90] (4,2);

    \draw[ms arrow](5.8,-.3)node[right]{Comp. G.}
        to[out=180,in=90] (3.9,-1);

    \draw[ms arrow](5.9,5)node[right]{Wind G.}
        to[out=180,in=90] (3.6,5);

    \draw[ms arrow](5.9,8.4)node[right]{Friction G.}
        to[out=180,in=90] (3.2,8);

    \draw[ms arrow,msColor1](5.3,-4.2)node[right]{G. Cell}
        to[out=180,in=90] (0,-2.5);

    \draw[ms arrow,msColor1](4.3,-1.9)node[right]{Nested G.}
        to[out=180,in=90] (2,-.5);

    \draw[ms arrow](4,-6)node[right]{Bathymetry}
        to[out=180,in=90] (2,-5);
    %drawing points on grid's conrners.
    \fill[msInk,ms body]
        (-5,-4.3) node [above] {$P_{1}$}
        (-.3,-5.9) node [below] {$P_{2}$}
        (5.5,-4) node [above] {$P_{3}$};
\end{tikzpicture}

\end{document}
